% Pursuing Stacks (cotes 134-2 to 134-8): archive cote 134-3, site page 64 (PDF page 65), typescript page 64.
% typescript. Draft compared: ps2.tex:3299-3369.
% AI-transcribed by Claude Opus 5.5 (claude-opus-5-5, adaptive thinking, effort high), Message Batches API, batch ps1 (msgbatch_01Rv5hCaNTtUt4JHx1yqoHPL), 2026-11-12; image review recorded in records/CHECK_DECISIONS_134_3_4.tsv.
\scanpage{134-3/64}
\mspage{\hand{64}}

theorem in the present special case can be formulated thus:

\uline{Corollary} 1  Assume the Lawvere element $L_A$ in $A^{\hand{\wedge}}$ is aspheric over the fi\xout{na}%
nal object $e_{A^{\hand{\wedge}}}$ of $A^{\hand{\wedge}}$,  and that moreover the latter b\uline{e} aspheric i.e. A
aspheric.  Then A is a test category in the wide sense\deleted{\illegible{1}}\hand{,} namely for any
C in (Cat),  the canonical functor  $\hand{i}_A\hand{j}_A(C)\hand{\longrightarrow} C$  is a weak equivalence
(and even aspheric\deleted{\xout{xx orxequivalently}} -  see prop. \hand{in} \deleted{p.}\inserted{section} 30, \hand{(}and also
\deleted{\illegible{1}}\inserted{section 29)}, for equivalent formulations\deleted{)}).

Thus we \uline{did} get after all a ``handy criterion'' sufficient to ensure
\margin{25.3. or rather, than the set of conditions T1 to T3}this basic test-property,  which looks a lot less strong \deleted{than} a priori
than the condition (T 2) \inserted{[plus T1, T3]} of total asphericity of $A^{\hand{\wedge}}$.

But let's now come back to the more general situation of the theorem,
with a functor
\[ i:\ A\hand{\longrightarrow}(\mathrm{Cat}) \]
subjected only to the mild condition that the categories $\hand{i}(\hand{a})$ (for a in
A) have final objects.  Assume condition b (iii) to be satisfied, namely
that $i^{\hand{*}}(\hand{\Delta}_1)$ is aspheric over $e_{A^{\hand{\wedge}}}$,  and therefore it is a homotopy inter\xout{\unsure{va}}%
val in $A^{\hand{\wedge}}$.  \over{W}{L}et again $L_A$ be the Lawvere element in $A^{\hand{\wedge}}$, we define (inde-%
pendently of any assumption on i ) a morphism of ``intervals'' in $A^{\hand{\wedge}}$,  com-%
patible even with the natural composition laws (by intersection) on both
members
\[ \hand{\phi} : J = i^{\hand{*}}(\hand{\Delta}_1)\hand{\longrightarrow} L_A \quad . \]
For this we remember that  for a in J , we get
\[ J(a)\hand{\simeq} \mathrm{Crib}\ i(a) \ \hand{\ensuremath{\hookrightarrow}}\ \text{set of all subobjects of } i(a) \text{ in } (\mathrm{Cat} \]
% NOTE: the end of this line ("(Cat") is cut at the right edge of the scan
(this bijection and the inclusion being functorial in a),  thus if C is
in J(a)  i.e. a crible in i(a),  we associate to this
\[ \hand{\phi}(C) = \text{ subobject of a in } A^{\hand{\wedge}}, \text{ corresponding to the crible in} \]
\[ \text{A/a of all b}\hand{/}\text{a such that } i(b)\hand{\to} i(a) \text{ factors through} \]
\[ \text{the crible } C\hand{\subset} i(a) \quad ; \]
it is immediate that the map  $\hand{\phi}_a : J(a)\hand{\longrightarrow} L(a)$  thus obtained is functo-%
rial in a for variable a,  hence a \xout{morphism} map $\hand{\phi} : J\hand{\longrightarrow} L$ ,  and it is
immediately checked too that this is ``compatible with endpoints'' - namely
when C is full respectively empty,  then so is $\hand{\phi}(C)$  (for the ``empty''
case,  this comes from the fact that the categories i(a) are non empty).
Applying now the comparison lemma for homotopy intervals (p. 60),  we get
the following

\uline{Corollary} 2:  Under the general conditions of the theorem,  and assuming
moreover that $i^{\hand{*}}(\hand{\Delta}_1)$ is aspheric over $e_{A\overset{\hand{\wedge}}{u}}$ (condition b (iii)),  it fol\xout{l}-%
lows that the Lawvere element $L_A$ of $A^{\hand{\wedge}}$is aspheric over $e_{A^{\hand{\wedge}}}$ .  Assume
moreover that $e_{A^{\hand{\wedge}}}$ is aspheric i.e. A aspheric.  T\over{h}{\unsure{h}}en we get :
